\section{Task construction and experimental design}
\label{sec:design}
\paragraph{Source validity and changed assumptions.}
Each family pairs a source setting where a procedure is valid with a target feature request that changes a relevant assumption. The final suite contains six F-series and seven X-series synthetic families, covering parsing, lifetimes, paths, rendering, identity, confidentiality, and state consistency. Existing baseline code is visible in every condition. N removes supplied memory, not the task's implementation. Appendix~\ref{app:families} gives the family contracts and construction records.

\paragraph{Reference controls.}
The saved preflight checks three target states in each family: the incomplete baseline must fail functionality and pass the focal witness; an unadapted functional reference must pass functionality and fail the witness; and a repaired reference must pass both. Thus avoiding the requested feature is insufficient for success. The final manifest fixes the synthetic artifacts and source procedures. Constructors and evaluated agents are separate roles; constructors do not receive sealed witnesses or evaluated agent outcomes. The controls establish these three states, while validation of alternative implementations requires examining what the observer actually measures.

\begin{table}[t]
\centering\small\begin{tabular}{lllr}
\toprule
Configuration & Agent / tools & Sampling selection & Recorded / valid\\
\midrule
GPT-5.5 Low & Codex / shell, patch & low reasoning & 104 / 96\\
Luna Medium & Codex / shell, patch & medium reasoning & 104 / 97\\
Terra Medium & Codex / shell, patch & medium reasoning & 104 / 94\\
Qwen Next & MiniSWE / bash & $T=1$ & 104 / 104\\
Qwen 30B & MiniSWE / bash & $T=0.7$ & 104 / 102\\
Devstral & MiniSWE / bash & $T=0.15$ & 104 / 102\\
\bottomrule
\end{tabular}

\caption{Executed design: thirteen families, four conditions, and two repetitions per configuration. Codex uses 480 seconds and 32 command items; MiniSWE uses 900 seconds and 30 steps, a 32,768-token context, and 512 output tokens per request. Appendix~\ref{app:config} records exact requested IDs, revisions, settings, and tool details.}
\label{tab:design}
\end{table}

\paragraph{Configurations and execution.}
The six configurations each complete their 104-run schedule, producing \RecordedRuns{} recorded runs: \ValidRuns{} technically valid and \InvalidRuns{} invalid. Codex and MiniSWE have different native interfaces, instruction envelopes, and budgets. Requested aliases are preserved; backend identities are not inferred. Executed manifests record run order. A seventh planned configuration is excluded because quota exhaustion left its cohort unfinished; this retrospective scope decision is documented, and none of its outcomes or examples are used (Appendix~\ref{app:corrections}).

\paragraph{Memory conditions.}
The conditions are no memory (N), correct source memory (C), matched irrelevant memory (I), and correct source memory with a generic applicability boundary (B). C, I, and B each contain a 4,096 byte envelope including neutral padding; N has none. Bytes, not tokens, are equalized. The B reminder asks whether the source assumption still holds and supplies no target repair. C$-$N is the frozen primary contrast and estimates the effect of supplying the correct source memory as a whole, including any generic effect of added context. I$-$N and C$-$I help diagnose that contribution but do not perfectly isolate semantic memory content. B$-$C tests the generic boundary reminder.

\paragraph{Outcomes and scoring.}
For a technically valid run, $F$ denotes functionality pass and $S$ focal security witness pass. The primary endpoint is
\begin{equation} U=F\land\neg S. \end{equation}
All four joint outcomes are reported. In particular, $\neg F\land S$ is a functional failure, while $F\land S$ means \emph{functionality and focal witness pass}. Technical invalids retain missing F/S/U values in the primary table. Termination labels do not override evaluations.

A versioned validity review changed four Devstral focal witness labels because the X05 and X28 observers had inspected optional internal ledgers rather than required return values and effects. The correction is narrowly scoped, preserves the original scores, changes no functionality or validity label, and leaves the qualitative C$-$N conclusion unchanged. Appendix~\ref{app:corrections} gives the full scoring history. X06 retains its executed confidentiality scope: rejecting the mixed record can pass without demonstrating retention of useful public fields.

\paragraph{Primary analysis.}
A family is eligible for a contrast when both repetitions in both compared conditions are valid with observed F and S. We average repetitions within family and condition, then average paired differences equally over eligible families. Tables report the eligible count. Exploratory 95\% percentile intervals use 20,000 paired family bootstrap resamples. Exact sign flip and family omission diagnostics appear in Appendix~\ref{app:sensitivity}. These are comparisons within configurations, not a ranking of model security.

\paragraph{Measurement availability and bounds.}
A separate review distinguishes interruption, timeout, candidate execution exception, and evaluator incompatibility. Completed outcome components are retained, while unresolved components remain unknown. Appendix~\ref{app:missingness} reports treatment specific counts and bounds over every assignment of the unresolved outcomes consistent with the completed observations.
